\documentclass[12pt,a4paper]{amsart}

\usepackage{amsmath,amssymb}
\usepackage{enumerate}
% \usepackage[english,brazil]{babel}
\usepackage[portuguese]{babel}
\usepackage{ragged2e}
\usepackage{blindtext}
\usepackage{amsfonts}
\usepackage{amsmath}
\usepackage{mathtools}
\usepackage[dvipsnames]{xcolor}

\definecolor{bege}{HTML}{FDFFD1}
\pagecolor{bege}


\begin{document}

\title{Simulação numérica do problema de N-corpos gravitacional}
\author{
    Octavio Augusto Potalej
    \vspace{0.25cm} \\ 
    IME-USP
}

\maketitle

\vspace{1cm}
\justifying

\begin{center}
    \begin{minipage}{10cm}
    \Large{
    
    O Problema de N-corpos gravitacional consiste em determinar a dinâmica de $N$ corpos celestes que interagem gravitacionalmente. Em geral, o problema não admite solução analítica fechada, então simulações numéricas tornam-se uma ferramenta útil para sua investigação.
    
    Nesta apresentação, veremos algumas abordagens numéricas para o problema, explorando diferentes métodos de integração, com foco especial em integradores simpléticos e na preservação da estrutura Hamiltoniana do sistema.

    % Resumo com, no máximo, 14 linhas (não mudar o tamanho da fonte e nem a formatação do texto)
    
    }
    \end{minipage}
 \end{center}
	
\vspace{1.5cm}

 \begin{thebibliography}{99}

    \bibitem[{\bf 1}]{1} S. J. Aarseth, {\em Gravitational N-Body Simulations: Tools and Algorithms}, Cambridge University Press, 2003.
    \bibitem[{\bf 2}]{2} E. Hairer; C. Lubich; G. Wanner, {\em Geometric Numerical Integration}, Springer, 2006.
        
\end{thebibliography}

\bigskip

\begin{flushleft}
    \emph{Email:}  \texttt{oapotalej@ime.usp.br}
\end{flushleft}	
	
\end{document}
